
\documentclass{article}
\usepackage{amsmath}
\usepackage[top=3cm, bottom=3cm, outer=3cm, inner=3cm]{geometry}
\usepackage{multicol}
\usepackage{graphicx}
\usepackage{booktabs}
\usepackage{pdflscape}
\usepackage{url}
\usepackage{hyperref}
\usepackage{array}
\newcolumntype{x}[1]{>{\centering\arraybackslash\hspace{0pt}}p{#1}}
\usepackage[natbibapa]{apacite}
\usepackage{pdfpages}
\usepackage{multirow}
\usepackage[normalem]{ulem}
\useunder{\uline}{\ul}{}
\usepackage{svg}
\usepackage[T1]{fontenc}
\usepackage{xcolor}
\usepackage{listings}
\lstdefinestyle{ascii-tree}{
    literate={├}{|}1 {─}{--}1 {└}{+}1 
  }
\lstset{basicstyle=\ttfamily,
  showstringspaces=false,
  commentstyle=\color{red},
  keywordstyle=\color{blue}
}
%\usepackage{booktabs}
\usepackage{caption}
\usepackage{subcaption}
\usepackage{float}
\newcolumntype{M}[1]{>{\centering\arraybackslash}m{#1}}
\newcolumntype{N}{@{}m{0pt}@{}}


%%%%%%%%%%%%%%%%%%%%%%%%%%%%%%%%%%%%%%%%%%%%%%%%%%%%%%%%%%%%%%%%%%%%%%%%%%%%
%%%%%%%%%%%%%%%%%%%%%%%%%%%%%%%%%%%%%%%%%%%%%%%%%%%%%%%%%%%%%%%%%%%%%%%%%%%%
\newcommand{\itemEmail}{ccapatinta@unsa.edu.pe}
\newcommand{\itemStudent}{Camilo B.A. Capatinta Almanza}
\newcommand{\itemCourse}{Organización y Métodos}
\newcommand{\itemCourseCode}{2502227}
\newcommand{\itemSemester}{IV}
\newcommand{\itemUniversity}{Universidad Nacional de San Agustín de Arequipa}
\newcommand{\itemFaculty}{Facultad de Ingeniería de Producción y Servicios}
\newcommand{\itemDepartment}{Departamento Académico de Ingeniería de Sistemas e Informática}
\newcommand{\itemSchool}{Escuela Profesional de Ingeniería de Sistemas}
\newcommand{\itemAcademic}{2026 - B}
\newcommand{\itemInput}{02/10/2026}
\newcommand{\itemOutput}{05/10/2026}
\newcommand{\itemDocente}{William M. Bornas Rios}
\newcommand{\itemTheme}{Minería de procesos, ERP, ROI y Kardex}
\newcommand{\itemPracticeNumber}{}
%%%%%%%%%%%%%%%%%%%%%%%%%%%%%%%%%%%%%%%%%%%%%%%%%%%%%%%%%%%%%%%%%%%%%%%%%%%%
%%%%%%%%%%%%%%%%%%%%%%%%%%%%%%%%%%%%%%%%%%%%%%%%%%%%%%%%%%%%%%%%%%%%%%%%%%%%

\usepackage[english,spanish]{babel}
\usepackage[utf8]{inputenc}
\AtBeginDocument{\selectlanguage{spanish}}
\renewcommand{\figurename}{Figura}
\renewcommand{\refname}{Referencias}
\renewcommand{\tablename}{Tabla} %esto no funciona cuando se usa babel
\AtBeginDocument{%
	\renewcommand\tablename{Tabla}
}

\usepackage{fancyhdr}
\pagestyle{fancy}
\fancyhf{}
\setlength{\headheight}{40.51407pt}
\renewcommand{\headrulewidth}{1pt}
\renewcommand{\footrulewidth}{1pt}
\fancyhead[L]{\textbf{\itemUniversity}}
\fancyhead[C]{\fontsize{7}{7}\selectfont\itemFaculty \\ \itemDepartment \\ \itemSchool \\ \textbf{\itemCourse}}
\fancyhead[R]{\textbf{\itemAcademic}}
\fancyfoot[L]{Estudiante Capatinta Camilo}
\fancyfoot[C]{\itemCourse}
\fancyfoot[R]{Página \thepage}

% Para el código fuente
\usepackage{color, colortbl}
\definecolor{dkgreen}{rgb}{0,0.6,0}
\definecolor{gray}{rgb}{0.5,0.5,0.5}
\definecolor{mauve}{rgb}{0.58,0,0.82}
\definecolor{codebackground}{rgb}{0.95, 0.95, 0.92}
\definecolor{tablebackground}{rgb}{0.8, 0, 0}

\lstset{frame=tb,
	language=bash,
	aboveskip=3mm,
	belowskip=3mm,
	showstringspaces=false,
	columns=flexible,
	basicstyle={\small\ttfamily},
	numbers=none,
	numberstyle=\tiny\color{gray},
	keywordstyle=\color{blue},
	commentstyle=\color{dkgreen},
	stringstyle=\color{mauve},
	breaklines=true,
	breakatwhitespace=true,
	tabsize=3,
	backgroundcolor= \color{codebackground},
}

% Título del índice en español y enlaces en el índice
\addto\captionsspanish{\renewcommand{\contentsname}{Índice}}
\hypersetup{
    colorlinks=true,
    linkcolor=black,   % enlaces del índice
    citecolor=black,
    urlcolor=black,
    linktoc=all        % toda la línea del índice es clicable
}

% Evita que natbib imprima su propio título "Referencias"
% (usamos nuestra propia \section para que aparezca numerada en el índice)
\makeatletter
\renewcommand{\bibsection}{}
\makeatother

\begin{document}
	
	\vspace*{10px}
	
	\begin{center}	
		\fontsize{17}{17} \textbf{ Ensayo }
	\end{center}
	\centerline{\textbf{\Large Tema: \itemTheme}}
	%\vspace*{0.5cm}	

	\begin{flushright}
		\begin{tabular}{|M{2.5cm}|N|}
			\hline 
			\rowcolor{tablebackground}
			\color{white} \textbf{Nota}  \\
			\hline 
			     \\[30pt]
			\hline 			
		\end{tabular}
	\end{flushright}	

	\begin{table}[h!]
		\begin{tabular}{|x{4.7cm}|x{4.8cm}|x{4.8cm}|}        
			\hline 
			\rowcolor{tablebackground}
			\color{white} \textbf{Estudiante} & \color{white}\textbf{Escuela}  & \color{white}\textbf{Asignatura}   \\
			\hline 
			{\itemStudent \par \itemEmail} & \itemSchool & {\itemCourse \par Semestre: \itemSemester \par Código: \itemCourseCode}     \\
			\hline 			
		\end{tabular}
	\end{table}		
	
	\begin{table}[H]
		\begin{tabular}{|x{4.7cm}|x{4.8cm}|x{4.8cm}|}
			\hline 
			\rowcolor{tablebackground}
			\color{white}\textbf{Docente} & \color{white}\textbf{Tema}  & \color{white}\textbf{Duración}   \\
			\hline 
			\itemDocente & \itemTheme & 04 horas   \\
			\hline 
		\end{tabular}
	\end{table}
	
	\begin{table}[H]
		\begin{tabular}{|x{4.7cm}|x{4.8cm}|x{4.8cm}|}
			\hline 
			\rowcolor{tablebackground}
			\color{white}\textbf{Semestre académico} & \color{white}\textbf{Fecha de inicio}  & \color{white}\textbf{Fecha de entrega}   \\
			\hline 
			\itemAcademic & \itemInput &  \itemOutput  \\
			\hline 
		\end{tabular}
	\end{table}
    \clearpage
    \tableofcontents
    \clearpage

    \section{Actividades}

        Para los algoritmos de Bor\r{u}vka, Jarník, Kruskal, Dijkstra y Floyd
        realice una tabla comparativa de tiempos de ejecución para 10 nodos,
        analizando:

        \begin{itemize}
        \item Complejidad temporal
        \item Complejidad espacial
        \item Casos de uso, cuándo usar y cuándo no
        \end{itemize}

    \subsection{Comparación de algoritmos}

    \begin{table}[htbp]
        \centering
        \caption{Complejidad temporal y espacial para grafos con $n=10$ nodos}
        \label{tab:complejidades}
        \begin{tabular}{l c c c}
            \toprule
            Algoritmo & Tiempo de ejecución ($n=10$) & Complejidad temporal & Complejidad espacial \\
            \midrule
            Borůvka \cite{boruvka1926} & $\approx 149$ operaciones (grafo denso, $E\log V$) & $O(E\log V)$ & $O(V+E)$ \\
            Jarník (Prim) \cite{jarnik1930} & $100$ operaciones con matriz; $\approx 149$ con cola de prioridad & $O(V^2)$ con matriz; $O(E\log V)$ con cola de prioridad; $O(E+V\log V)$ con montículo de Fibonacci & $O(V^2)$ con matriz; $O(V+E)$ con lista de adyacencia \\
            Kruskal \cite{kruskal1956} & $\approx 247$ operaciones ($E\log E$) & $O(E\log E) = O(E\log V)$ & $O(V+E)$ (lista de aristas + unión-búsqueda) \\
            Dijkstra \cite{dijkstra1959} & $100$ operaciones con matriz; $\approx 183$ con cola de prioridad & $O(V^2)$ con matriz; $O((V+E)\log V)$ con montículo binario; $O(E+V\log V)$ con montículo de Fibonacci & $O(V^2)$ con matriz; $O(V+E)$ con lista de adyacencia \\
            Floyd (Floyd--Warshall) \cite{floyd1962} & $1000$ operaciones ($V^3$) & $O(V^3)$ (fija, sin variantes relevantes) & $O(V^2)$ (matriz de distancias) \\
            \bottomrule
        \end{tabular}
    \end{table}

    \clearpage
    \subsection{Casos de uso}
    \begin{landscape}
    \begin{table}[htbp]
        \centering
        \caption{Casos de uso: cuándo usar y cuándo no usar cada algoritmo}
        \label{tab:casos-uso}
        \begin{tabular}{l p{7cm} p{7cm}}
            \toprule
            Algoritmo & Cuándo usar & Cuándo no usar \\
            \midrule
            Borůvka \cite{boruvka1926} & Construir un árbol de expansión mínima (MST) en grafos densos o en entornos paralelos/distribuidos, ya que procesa todos los componentes simultáneamente en cada ronda; es la base de algoritmos paralelos modernos de MST. & Si solo se necesita el MST de un grafo pequeño o disperso, Kruskal o Prim son más simples de implementar y depurar; tampoco aplica si el grafo tiene pesos negativos relevantes para otro objetivo (el MST no depende del signo, pero el algoritmo no resuelve caminos mínimos). \\
            Jarník (Prim) \cite{jarnik1930} & Construir el MST de grafos densos representados con matriz de adyacencia, o cuando conviene crecer el árbol desde un único vértice (por ejemplo, redes donde se parte de un nodo central). & En grafos dispersos representados con listas de aristas, Kruskal suele ser más eficiente y más simple; tampoco sirve si se requieren caminos mínimos (solo da un árbol de costo mínimo, no las distancias más cortas entre pares). \\
            Kruskal \cite{kruskal1956} & Construir el MST de grafos dispersos, o cuando las aristas ya están ordenadas o es fácil ordenarlas; es sencillo de combinar con una estructura de conjuntos disjuntos (unión-búsqueda). & En grafos muy densos, el costo de ordenar $E\log E$ aristas puede superar al de Prim con matriz ($O(V^2)$); tampoco aplica si el grafo no es conexo y se requiere un único árbol (produce un bosque de expansión mínima en ese caso). \\
            Dijkstra \cite{dijkstra1959} & Calcular el camino más corto desde un único origen hacia todos los demás nodos, en grafos con pesos no negativos; es la opción estándar para ruteo en redes y mapas. & Si el grafo tiene aristas de peso negativo, el algoritmo puede dar resultados incorrectos (no está diseñado para eso); tampoco conviene si se necesitan los caminos más cortos entre todos los pares de nodos y el grafo es denso, donde Floyd–Warshall resulta más directo. \\
            Floyd (Floyd–Warshall) \cite{floyd1962} & Calcular los caminos más cortos entre todos los pares de nodos a la vez, incluso con pesos negativos (siempre que no haya ciclos negativos); útil en grafos pequeños o densos donde $O(V^3)$ es aceptable. & En grafos grandes o dispersos, el costo $O(V^3)$ resulta excesivo comparado con ejecutar Dijkstra desde cada nodo ($O(V\cdot(V+E)\log V)$); tampoco es la opción natural si solo se necesita el camino más corto desde un único origen. \\
            \bottomrule
        \end{tabular}
    \end{table}
    \end{landscape}

    \section{Referencias bibliográficas}
    \bibliographystyle{apacite}
    \bibliography{referencias}

    \end{document}